\chapter{Resultados}
\label{cap:resultados}

\section{Análise comparativa das especificações 0126 e 1580}

\subsection{Campo de aplicação das especificações}
\label{subsec:resultado-abrangencia}

A comparação da seção~1 das especificações evidencia uma ampliação do campo de aplicação da 1580 em relação à 0126 (requisito RN-0001). Enquanto a 0126 delimitava sua aplicação a empreendimentos constituídos por um único prédio e remetia os conjuntos de prédios ou blocos à especificação 0128, a 1580 inclui expressamente os empreendimentos de conexão coletiva horizontais e verticais. O item~7.8.1 da 1580 complementa essa delimitação ao definir a situação em que unidades contíguas ou geminadas, simultaneamente sem área de uso comum e alinhadas à via pública, devem ser atendidas de forma individual \cite[seção~1, p.~7]{enel0126}\cite[seção~1, p.~3, e item~7.8.1, p.~25]{enel1580}.

A diferença é classificada como alterado, conforme os critérios do Quadro~\ref{qua:classificacao-normativa}, por modificar o alcance de aplicação da especificação de conexão coletiva em substituição à 0126. A ampliação não representa o início da regulamentação dos condomínios horizontais no Ceará, uma vez que a própria 0126 previa a especificação 0128 como referência para conjuntos de edificações. O histórico de versões da 1580 identifica a substituição da 0126, mas não relaciona expressamente a 0128 entre os documentos cancelados; portanto, a comparação não permite concluir pela revogação integral desta última \cite[seção~1]{enel0126}\cite[item~2, p.~3--4]{enel1580}.

A configuração residencial vertical do empreendimento, apresentada na Seção~\ref{sec:caracterizacao-empreendimento}, permanece compreendida no campo de aplicação de ambas as especificações. A ampliação do campo de aplicação da 1580, isoladamente, não implica mudança na configuração de atendimento do empreendimento. Por essa razão, o requisito RN-0001 integra a comparação normativa geral, mas não é selecionado para a avaliação dos impactos técnicos, conforme os critérios da Seção~\ref{sec:impactos}.

\subsection{Enquadramento de instalações de uso coletivo no Grupo B}
\label{subsec:resultado-grupob-coletivo}

O art.~23, §~2º, da Resolução Normativa ANEEL nº~1.000/2021 admite o enquadramento no Grupo B de instalações de uso coletivo com carga superior a 75~kW em edificações de múltiplas unidades consumidoras, condicionado à predominância de unidades do Grupo B, superior à metade do total, à solicitação ou concordância do consumidor e à avaliação de viabilidade pela distribuidora mediante estudo \cite[art.~23, §~2º]{aneel1000}.

A especificação 0126 reproduz essas condições no item~7.2.1.5 e identifica, em seu histórico de versões, a adequação à REN 1000 como uma das mudanças da versão 03. Na 1580, as três condições gerais constam da nota do item~7.8.2; a Nota~5 da seção~7.9 também as menciona em relação a novas instalações de uso coletivo \cite[seção~2, p.~8, e item~7.2.1.5, p.~16]{enel0126}\cite[item~7.8.2, p.~25, e seção~7.9, Nota~5, p.~28]{enel1580}.

A possibilidade de enquadramento tem origem regulatória anterior à 1580 e já estava prevista na 0126. A transferência da prescrição de um subitem próprio para notas na estrutura da 1580 motiva sua classificação como reorganizado (RN-0002), preservado o núcleo das três condições regulamentares. A expressão relativa às unidades que já pertençam ao Grupo B, empregada na 1580, deve ser interpretada à luz do enquadramento previsto na resolução, sem inferir uma exigência adicional de que as unidades estejam previamente energizadas. A classificação restringe-se à regra geral de enquadramento; não estabelece equivalência entre limites técnicos, requisitos de proteção, procedimentos documentais ou condições regionais das duas especificações.

O memorial de cálculo elétrico do empreendimento registra 95,713~kW de carga instalada nas instalações comuns do condomínio, valor distinto da carga instalada total do empreendimento \cite[p.~6]{palatiummemorial2025}. Como a carga instalada do condomínio supera 75~kW, a regra geral de enquadramento é pertinente ao estudo de caso. Entretanto, não é selecionada isoladamente para a avaliação dos impactos, pois não foi demonstrada alteração técnica dessa regra em relação à 0126. A análise não confirma o atendimento a todas as condições de enquadramento nem substitui o estudo de viabilidade da distribuidora.

\subsection{Critérios de cálculo de demanda}
\label{subsec:resultado-demanda}

A seção~7.11 da especificação 0126 estabelece os métodos da área útil e da carga instalada, cujas expressões estão apresentadas nas Equações~\ref{eq:demanda-area} a~\ref{eq:demanda-carga-instalada-0126}. Os itens~7.11.1 a~7.11.3, as Tabelas~1, 2, 4, 5, 6, 15 e~16 e os exemplos do Anexo G permitem identificar, no próprio documento, os dados necessários, os coeficientes de cálculo, os critérios de composição da demanda e as condições de aplicação. No método da área útil, o dispositivo numerado como 7.12.2.9 na página~40 estabelece margem mínima de segurança de 40~\% sobre a demanda dos apartamentos e requer justificativa para um valor superior. A possibilidade de propor outro método no contexto do critério da carga instalada consta do item~7.11.3.5. A determinação da demanda já era de responsabilidade do responsável técnico pelo projeto, conforme o item~7.11.1.5 \cite[seção~7.11, p.~38--41, e Anexo G]{enel0126}.

Na especificação 1580, a alínea~k do item~7.3.2 (p.~21) exige o memorial detalhado da demanda da entrada geral, dos ramais de distribuição principal, dos ramais alimentadores das unidades, dos barramentos blindados e dos centros de medição, em kVA. A nota ao final desse item (p.~22) atribui a determinação da demanda prevista ao responsável técnico, empregando a expressão \emph{utilizando de metodologias recomendadas pelas distribuidoras em documentos específicos}. A nota após o item~7.11.9 (p.~43) reafirma a responsabilidade, mas emprega a expressão \emph{podendo ser utilizada metodologias recomendadas pelas distribuidoras em documentos específicos} \cite[item~7.3.2 e nota após o item~7.11.9]{enel1580}. A primeira redação aponta para o emprego de metodologias recomendadas; a segunda apresenta esse emprego em termos de possibilidade. Nenhuma das duas notas identifica os títulos ou códigos desses documentos. Essa diferença de redação, isoladamente, não permite concluir que haja liberdade irrestrita de escolha do método nem que determinada metodologia seja obrigatória ou proibida.

O Quadro~\ref{qua:demanda-comparacao} sistematiza as diferenças entre o tratamento metodológico, a responsabilidade técnica e a documentação dos cálculos.

\begin{quadro}[htbp]
  \centering
  \begin{Spacing}{1.}
    \caption{Comparação do tratamento do cálculo de demanda}
    \label{qua:demanda-comparacao}
  \end{Spacing}
  {\small
  \renewcommand{\arraystretch}{1.2}
  \setlength{\tabcolsep}{4pt}
  \begin{tabular}{|>{\raggedright\arraybackslash}p{0.21\linewidth}|>{\raggedright\arraybackslash}p{0.33\linewidth}|>{\raggedright\arraybackslash}p{0.34\linewidth}|}
    \hline
    \textbf{Aspecto} & \textbf{Especificação 0126} & \textbf{Especificação 1580} \\\hline
    Método geral (RN-0003) & Seção 7.11, tabelas e Anexo G: procedimentos, equações e exemplos. & Remissões genéricas a metodologias recomendadas nas notas das p. 22 e 43; não reproduz o procedimento geral da 0126. \\\hline
    Responsabilidade (RN-0004) & Item 7.11.1.5: atribuição ao responsável técnico pelo projeto. & Notas das p. 22 e 43: preservação da atribuição ao responsável técnico. \\\hline
    Memorial (RN-0005) & Item 7.12.2.3, alínea c: cálculo de demanda e justificativa do dimensionamento. & Item 7.3.2, alínea k: discriminação dos elementos cujo cálculo detalhado da demanda deve ser apresentado. \\\hline
  \end{tabular}}
  \Fonte{Elaborado pelo autor com base nas especificações 0126 e 1580 \cite[seções~7.11 e~7.12, Anexo G]{enel0126}\cite[p.~21--22 e~43]{enel1580}.}
\end{quadro}

O tratamento do método geral é classificado como alterado (RN-0003) pela redução do detalhamento e pelas remissões a metodologias recomendadas na 1580. A atribuição ao responsável técnico é classificada como reorganizada (RN-0004), pois a responsabilidade prevista na 0126 passa a constar em notas. A discriminação dos elementos cujo cálculo deve integrar o memorial é classificada como alterada (RN-0005). A 1580 também contém critérios específicos de demanda, como o tratamento dos sistemas de incêndio e a seleção de condutores e proteções em função da corrente de demanda e da queda de tensão \cite[itens~7.11.9 e~7.19.6]{enel1580}.

A consulta descrita na Seção~\ref{sec:procedimento} esclareceu que as notas da 1580 não identificam um documento complementar específico para o Ceará. Segundo a resposta do setor Loja Clientes Corporativos (CE), recebida em 23 de novembro de 2026, os critérios do item~7.11 da 0126 perderam a vigência normativa obrigatória e podem ser usados como referência técnica, em conjunto com os critérios da ABNT NBR 5410, desde que a metodologia e os fatores adotados sejam explicitados na memória de cálculo. A margem mínima de 40~\% deixa, assim, de ser exigível por força da 0126. A determinação da demanda permanece sob responsabilidade do responsável técnico, com base nas metodologias aplicáveis.

\subsection{Limites de atendimento e alternativas de transformação (REVISAR)}
\label{subsec:resultado-limites}

O limite geral de 75~kW de carga instalada por unidade consumidora em baixa tensão, apresentado nas seções~7.2 e~7.3 da 0126, permanece na alínea~a do item~7.8.2 da 1580. A manutenção desse valor é classificada como reorganizado (RN-0006), considerando sua apresentação em um dispositivo específico para o atendimento individual. A possibilidade regulamentar de enquadramento das instalações de uso coletivo acima de 75~kW no Grupo B foi analisada no requisito RN-0002 e não constitui uma nova alteração. O Comunicado Técnico nº~02/2022, de 28 de março de 2022, da Enel Distribuição São Paulo, também registrou essa disposição do art.~23 da REN 1000, mas seu campo de aplicação restringe-se à área de concessão paulista \cite[item~7.2.1.5]{enel0126}\cite[item~7.8.2]{enel1580}\cite[seções~3 e~5]{enelct022022}. O limite paulista de 225~kVA previsto nas disposições específicas da 1580 não deve ser transposto indistintamente para novas conexões no Ceará.

Além do limite geral, a 0126 especifica condições técnicas para o atendimento individual em suas seções~7.2 e~7.3: proteção trifásica nominal de até 125~A e, para a unidade dos serviços comuns com carga acima de 75~kW, apresentação e aprovação do cálculo de demanda, proteção de até 125~A e separação física e elétrica dos circuitos de outros pontos de entrega (itens~7.2.1.2 e~7.3.1.2). Na 1580, a nota do item~7.8.2 estabelece as condições regulamentares do enquadramento; a Tabela~8 apresenta modalidades por tipo de rede, inclusive faixa de 26 a 100~kW para atendimento trifásico na rede subterrânea do Ceará, com asterisco. O item~7.17.1 admite, sob as condições nele indicadas, medição direta com corrente de demanda de 101 a 200~A. Esses parâmetros não são diretamente intercambiáveis: carga instalada, corrente de demanda do medidor e corrente nominal da proteção são grandezas e condições distintas \cite[itens~7.2.1 e~7.3.1]{enel0126}\cite[item~7.8.2, Tabela~8, p.~32, e item~7.17.1]{enel1580}. A correspondência integral das condições técnicas aplicáveis à unidade de serviços comuns acima de 75~kW é inconclusiva (RN-0009), sem presumir que o limite de 125~A tenha sido universalmente retirado ou que 200~A seja o novo limite de proteção.

Para o Ceará, a Tabela~8 da 1580 explicita faixas de carga para o fornecimento bifásico segundo o tipo de rede e o sistema de distribuição. A Tabela~8 da 0126 já previa critérios de poste auxiliar para ramais bifásicos. Esse antecedente sustenta a classificação como alterado (RN-0011) quanto à explicitação das condições de atendimento, preservando a distinção entre os critérios de fornecimento e os de dimensionamento do poste. A alteração não incide na configuração trifásica do Palatium \cite[tabela~8, p.~66]{enel0126}\cite[item~7.10.2, tabela~8]{enel1580}.

Para a entrada geral, o item~7.5.1.1 da 0126 admite novos PMUC de demanda total de até 300~kVA na rede externa de baixa tensão, desde que todas as unidades possam ser atendidas nessa tensão. O item~7.5.1.2 prevê eventual extensão ou reforço de rede e instalação ou substituição do transformador. A alínea~k do item~7.5.1.4 admite avaliação, mediante solicitação formal, para demanda acima de 300~kVA e abaixo de 600~kVA, utilizando as soluções ali especificadas. Para demanda superior a 600~kVA, o texto prevê a possibilidade de entradas individualizadas em baixa tensão, mediante aprovação e separação das instalações comuns, enquanto o item~7.5.2.1 determina, como regra, conexão à rede externa de média tensão acima de 600~kVA ou quando houver unidade atendível em média tensão. Os itens~7.5.3.1 e~7.5.3.2 disciplinam o aumento de demanda em PMUC existentes, incluindo dois transformadores de 225 ou 300~kVA e, superada a capacidade de dois transformadores de 300~kVA, entrada em média tensão com subestação abrigada \cite[itens~7.5.1 a~7.5.3, p.~18--21]{enel0126}.

Na 1580, o item~7.8.3 vincula a demanda máxima do empreendimento à capacidade dos alimentadores do padrão construtivo definido na seção~7.11 e recomenda consulta prévia à distribuidora. A Tabela~9 organiza os tipos de conexão e postos de transformação por demanda e área de concessão. O Quadro~\ref{qua:limites-ceara} apresenta os arranjos expressamente aplicáveis ao Ceará, sem incluir as câmaras transformadoras exclusivas de São Paulo.

\begin{quadro}[htbp]
  \centering
  \begin{Spacing}{1.}
    \caption{Limites e soluções de conexão coletiva aplicáveis ao Ceará na 1580}
    \label{qua:limites-ceara}
  \end{Spacing}
  {\small
  \renewcommand{\arraystretch}{1.2}
  \setlength{\tabcolsep}{4pt}
  \begin{tabular}{|>{\raggedright\arraybackslash}p{0.50\linewidth}|>{\raggedright\arraybackslash}p{0.19\linewidth}|>{\raggedright\arraybackslash}p{0.18\linewidth}|}
    \hline
    \textbf{Solução} & \textbf{Demanda (kVA)} & \textbf{Item} \\\hline
    Transformador no poste da distribuidora & Até 300 & 7.11.1 \\\hline
    Transformador em poste no limite da propriedade & Até 300 & 7.11.2 \\\hline
    Centro de transformação de superfície compacto (CTSC) & Até 630 & 7.11.3 \\\hline
    Transformador pedestal & Até 1\,000 & 7.11.5 \\\hline
    Centro de transformação (CTS) & A partir de 2\,000 & 7.11.6 \\\hline
    Subestação de entrada de energia interna ao empreendimento (SEE) & Superior a 300 & 7.12.1 \\\hline
  \end{tabular}}
  \Fonte{Elaborado pelo autor com base na Tabela~9 da especificação 1580, p.~33 \cite[tabela~9]{enel1580}.}
\end{quadro}

O quadro reproduz as faixas de demanda da Tabela~9, não os valores de potência nominal de cada transformador nem uma escala de atendimento automaticamente disponível. Por exemplo, o item~7.11.3 estabelece demanda de até 630~kVA para o CTSC radial; o item~7.11.5 prevê até dois transformadores em pedestal de 500~kVA, e a Tabela~9 condiciona essa solução à consulta de viabilidade. O CTS deve ser interpretado também segundo as configurações e potências instaladas da Tabela~10. A nota regional da Tabela~9 recomenda consulta prévia no Ceará e no Rio de Janeiro em razão da disponibilidade de materiais e soluções construtivas. As notas gerais condicionam os arranjos às condições técnicas, operacionais e de manutenção, admitindo avaliação de outras soluções de SEE \cite[tabelas~9 e~10, itens~7.11.1 a~7.11.6]{enel1580}.

O valor de 300~kVA continua previsto para o arranjo da entrada coletiva com transformador no poste da distribuidora, conforme item~7.11.1, com avaliação da disponibilidade da rede existente. A preservação desse valor de referência é classificada como reorganizado (RN-0007), sem equiparar automaticamente todas as condições de atendimento. A organização dos limites coletivos e alternativas de transformação por padrões construtivos, em substituição à estrutura das faixas e exceções da seção~7.5 da 0126, configura alteração (RN-0008). As opções de maior demanda da 1580 não permitem concluir que a capacidade de uma ligação à rede externa de baixa tensão tenha sido elevada de 300 para 630 ou 1\,000~kVA, pois os limites correspondem a arranjos de transformação distintos. Tampouco se estabelece que o limite geral tenha diminuído pelo fato de o item~7.11.1 indicar 300~kVA.

\subsection{Ramal de conexão aéreo e poste auxiliar}
\label{subsec:resultado-ramal-poste}

O Quadro~\ref{qua:ramal-poste} compara as condições do ramal aéreo derivado da rede secundária e do poste particular.

\par\noindent
\begin{minipage}{\linewidth}
  \captionof{quadro}{Comparação dos requisitos do ramal aéreo e do poste auxiliar}
  \label{qua:ramal-poste}
  \vspace{0.7\baselineskip}
  \centering
  {\small
  \renewcommand{\arraystretch}{1.2}
  \setlength{\tabcolsep}{4pt}
  \begin{tabular}{|>{\raggedright\arraybackslash}p{0.19\linewidth}|>{\raggedright\arraybackslash}p{0.33\linewidth}|>{\raggedright\arraybackslash}p{0.36\linewidth}|}
    \hline
    \textbf{Requisito} & \textbf{0126} & \textbf{1580} \\\hline
    Vão aéreo & Até 40~m do poste da rede ao ponto de entrega. & Até 30~m da rede secundária ao ponto de conexão, observada a ancoragem no Desenho~1. \\\hline
    Alturas mínimas & 5,5~m para trânsito de veículos, 4,5~m em entradas de veículos e 3,5~m onde circulam apenas pedestres. & Mesmos valores nas situações comparáveis do Desenho~1. \\\hline
    Comprimento e resistência do poste & Seleção por seção, vão e modalidade do ramal. A Tabela~8 indica poste 300/9 para seção superior a 95~mm². & Comprimentos de 6 ou 7,5~m conforme a posição relativa à rede. A Tabela~18 diferencia ramal único e ramais múltiplos e prevê condições para comprimentos e esforços fora dos padrões. \\\hline
    Engastamento & Um décimo do comprimento total mais 0,60~m. & Mesma expressão na Equação~1. \\\hline
    Circuitos do ramal aéreo & Acima de 200~kVA, o ramal de ligação fica a critério da Enel Ceará. & A Tabela~17 vincula os circuitos à potência do transformador. A distribuidora pode adotar outra configuração conforme disponibilidade e critérios técnicos. \\\hline
  \end{tabular}
  \Fonte{Elaborado pelo autor com base nos itens~7.5.1.4 e~7.6.1.3 e nas Tabelas~8 e~11 da 0126, e nos itens~7.16.2 a~7.16.5, nas Tabelas~17 e~18 e no Desenho~1 da 1580 \cite{enel0126,enel1580}.}}
\end{minipage}
\par

O vão, os critérios do poste e a definição dos circuitos foram classificados como alterados (RN-0012, RN-0014 e RN-0016). As alturas mínimas e a expressão de engastamento foram mantidas (RN-0013 e RN-0015).

A alínea~k do item~7.16.2 da 1580 exige termo que ateste a resistência mecânica para ramal de entrada com seção superior a 35~mm², acompanhado de especificações técnicas e do documento de responsabilidade técnica de projeto e execução. A 0126 prevê critérios de dimensionamento do poste e documentação geral de responsabilidade técnica, sem termo equivalente vinculado a esse limiar nos dispositivos confrontados. A exigência específica foi classificada como acrescentada (RN-0017) \cite[item~7.12.2.3 e tabela~8]{enel0126}\cite[item~7.16.2, alínea~k]{enel1580}.

\section{Estudo de caso: Edifício Palatium}

\subsection{Carga instalada e previsão de recarga veicular}
\label{subsec:impacto-carga-instalada}

O conceito de carga instalada é mantido (RN-0010) nas duas especificações e corresponde à soma das potências nominais dos equipamentos aptos a operar na unidade consumidora \cite[item~7.1.3, p.~12]{enel0126}\cite[seção~6, p.~10]{enel1580}. No Palatium, o memorial declara 707,31~kW no centro de proteção geral (CPG). Esse valor é reconstituído pela soma de 26 apartamentos de 23,240~kW, 95,713~kW do condomínio e 7,360~kW do circuito de incêndio. A soma, \(26\times23{,}240+95{,}713+7{,}360\), resulta em 707,313~kW \cite[p.~1, 5--6 e~12]{palatiummemorial2025}.

O dimensionamento associado à medição de cada apartamento inclui a recarga e indica 26,940~kW por unidade. A diferença de 3,70~kW por apartamento totaliza 96,20~kW de recarga prevista. Mantidas as parcelas do condomínio e de incêndio, a reconstituição \(26\times26{,}940+95{,}713+7{,}360\) resulta em 803,513~kW, cerca de 13,60~\% acima do total reconstituído sem recarga \cite[p.~12--13]{palatiummemorial2025}\cite[prancha~003]{palatiumpranchas2025}.

A diferença de 96,20~kW já está presente nos dados do projeto original. Para a comparação, adotam-se 707,31~kW como carga instalada declarada e 803,51~kW como carga instalada reconstituída com a previsão de recarga. A demanda declarada de 291,33~kVA permanece como referência do estudo de caso.

\subsection{Implicações dos critérios de cálculo de demanda (REVISAR)}
\label{subsec:impacto-demanda}

O memorial adota o critério da área útil para dimensionar o CPG. Para os 26 apartamentos de 163~\(\mathrm{m}^2\), utiliza 3,34~kVA por apartamento e fator de diversificação 21,06, valores que também constam das Tabelas~15 e~16 da 0126. Pela Equação~\ref{eq:demanda-apartamentos}, o produto \(3{,}34\times21{,}06\) resulta em aproximadamente 70,34~kVA de demanda residencial antes da aplicação do multiplicador \cite[p.~7]{palatiummemorial2025}\cite[tabelas~15 e~16]{enel0126}.

O multiplicador 3,2 aplicado à parcela residencial corresponde a uma margem de 220~\%. Somadas as parcelas do condomínio e de incêndio, a Equação~\ref{eq:demanda-palatium-documentada} reproduz a demanda declarada no memorial. Na sequência, o memorial registra 442,76~A de corrente de projeto e proteção geral de 500~A \cite[p.~7]{palatiummemorial2025}.

\begin{equation}
 D_{\mathrm{T}}=(70{,}34\times3{,}2+58{,}88+7{,}36)~\mathrm{kVA}
 \cong291{,}33~\mathrm{kVA}.
 \label{eq:demanda-palatium-documentada}
\end{equation}

Mantendo as parcelas de 58,88~kVA do condomínio e 7,36~kVA do incêndio, a aplicação demonstrativa do multiplicador histórico mínimo de 1,4 resulta em \(70{,}34\times1{,}4+58{,}88+7{,}36\), ou 164,72~kVA. A diferença de 126,61~kVA em relação ao valor declarado ilustra o efeito do multiplicador na mesma expressão. A margem superior à mínima histórica pode estar associada à previsão de recarga, cuja contribuição à demanda coletiva não foi discriminada no memorial. A memória de cálculo consultada não permite estabelecer a motivação do valor 3,2 \cite[p.~7 e~13]{palatiummemorial2025}.

A resposta da Enel, apresentada na seção de comparação normativa, esclarece em que condições o método da 0126 pode ser usado como referência. A determinação de uma demanda para o cenário da 1580 ainda depende dos critérios de cálculo aplicáveis e da justificativa dos fatores para o empreendimento. Assim, a consequência de RN-0003 sobre a demanda da entrada geral permanece classificada como informação documental insuficiente (IP-0001). A atribuição ao responsável técnico (RN-0004) permanece fora da avaliação de impactos, pois foi reorganizada.

\subsection{Discriminação dos cálculos de demanda no memorial}
\label{subsec:impacto-discriminacao-demanda}

A alínea~k do item~7.3.2 da 1580 discrimina os elementos cujas demandas devem ser apresentadas no memorial em quilovolt-ampères (kVA) \cite[item~7.3.2, alínea~k, p.~21]{enel1580}. O Quadro~\ref{qua:discriminacao-demanda} indica o que foi identificado para cada elemento nos documentos examinados. A demanda do quadro de distribuição do condomínio (QDC) também está calculada no memorial, que registra 58,884~kVA para os serviços comuns \cite[p.~8--9]{palatiummemorial2025}.

\par\noindent
\begin{minipage}{\linewidth}
  \captionof{quadro}{Documentação da demanda por elemento no Edifício Palatium}
  \label{qua:discriminacao-demanda}
  \vspace{0.7\baselineskip}
  \centering
  {\small
  \renewcommand{\arraystretch}{1.2}
  \begin{tabular}{|>{\raggedright\arraybackslash}p{0.25\linewidth}|>{\raggedright\arraybackslash}p{0.69\linewidth}|}
    \hline
    \textbf{Elemento} & \textbf{Registro no projeto examinado} \\\hline
    Entrada geral & Demanda de 291,33~kVA apresentada na página~7 do memorial. \\\hline
    Ramais de distribuição principal & Trechos identificados nos esquemas do projeto. Cálculo de demanda por trecho não discriminado no memorial. \\\hline
    Alimentadores das unidades consumidoras & Parcialmente documentados. O memorial informa potência ativa, tensão, fator de potência, dados de corrente, distância, condutores e proteções nas páginas~12 e~13. \\\hline
    Centro de medição (CM) & Centro único com 28 medições na configuração original. Cálculo de demanda específico do CM não discriminado no memorial. \\\hline
    Barramento blindado & Não aplicável à distribuição por cabos representada no projeto. O CM contém barramentos internos de distribuição. \\\hline
  \end{tabular}
  \Fonte{Elaborado pelo autor com base no memorial e nas pranchas do projeto \cite[p.~5, 7--9 e~12--13]{palatiummemorial2025}\cite[pranchas~003--005]{palatiumpranchas2025}.}}
\end{minipage}
\par

Nas páginas~12 e~13 do memorial, os valores da coluna denominada \emph{Demanda} são coerentes com correntes de 38,39~A e 44,50~A, calculadas para a alimentação trifásica de 380~V. Com o fator de potência (FP) de 0,92 empregado no memorial, a relação \(S=P/\mathrm{FP}\) permite reconstituir 25,26~kVA para 23,24~kW sem recarga e 29,28~kVA para 26,94~kW com recarga. A reconstituição corresponde à hipótese trifásica usada no cálculo e se limita aos valores aparentes derivados das potências ativas informadas. A avaliação por fase e a verificação dos condutores dependem das condições de distribuição das cargas \cite[p.~12--13]{palatiummemorial2025}.

Os diagramas mostram a derivação independente para o quadro de transferência automática do incêndio antes da proteção geral e ajudam a identificar quais trechos devem ter a demanda explicitada \cite[pranchas~003--005]{palatiumpranchas2025}. Para o cenário de aplicação da 1580, a necessidade de discriminar os cálculos dos elementos aplicáveis e explicitar a metodologia e os fatores adotados constitui uma adequação potencialmente necessária de natureza documental (IP-0002, vinculado a RN-0005). A classificação se baseia no memorial e nas pranchas examinados e admite a existência de cálculos complementares fora do acervo disponível.

\subsection{Implicações dos limites de atendimento (REVISAR)}
\label{subsec:impacto-limites}

A prancha~005 mostra o ramal de conexão aéreo seguido de ramal de entrada subterrâneo. Os apartamentos, com 26,94~kW por unidade incluindo a recarga, permanecem na faixa geral do atendimento trifásico. A unidade do condomínio possui 95,713~kW de carga instalada, demanda de 58,88~kVA, corrente calculada de 89,49~A e proteção de 125~A. A 1580 admite solicitação de atendimento em baixa tensão acima de 75~kW em rede aérea, sujeita a consulta prévia, disponibilidade técnica e condições regulamentares. A faixa até 100~kW da Tabela~8 corresponde à rede subterrânea. A faixa de 101 a 200~A do item~7.17.1 se refere à corrente de demanda em condições de medição direta. O enquadramento efetivo do condomínio e a correspondência entre as condições de proteção permanecem em aberto (RN-0002 e RN-0009) \cite[p.~6--9 e~13]{palatiummemorial2025}\cite[prancha~005]{palatiumpranchas2025}\cite[p.~32--33 e item~7.17.1]{enel1580}.

A demanda original de 291,33~kVA, discutida na Seção~\ref{subsec:impacto-demanda}, corresponde a 97,11~\% do patamar de 300~kVA. A diferença é de 8,67~kVA. Um aumento de aproximadamente 2,98~\% sobre a demanda declarada atingiria esse patamar. A disponibilidade de rede e a configuração de transformação dependem da avaliação da distribuidora \cite[p.~7]{palatiummemorial2025}\cite[itens~7.5.1.1--7.5.1.2]{enel0126}\cite[item~7.11.1 e tabela~9]{enel1580}. Os registros RN-0006 e RN-0007 preservam os valores de referência e RN-0008 descreve alternativas cujo efeito específico não foi demonstrado para a demanda original. A seleção de RN-0009 permanece indeterminada.

\subsection{Ramal aéreo e poste auxiliar no Palatium}
\label{subsec:impacto-ramal-poste}

O comprimento do vão aéreo não foi identificado nas pranchas. As cotas da infraestrutura enterrada descrevem outro trecho da ligação. Assim, a repercussão da redução do vão máximo foi classificada como informação documental insuficiente (IP-0003) \cite[pranchas~001 e~005]{palatiumpranchas2025}.

O projeto indica poste auxiliar 300/9 e um ramal de entrada geral com dois conjuntos de condutores em paralelo, cada um com três fases de 150~mm² e neutro de 70~mm², ligados à mesma proteção geral de 500~A. A nota~4 da Tabela~18 da 1580 distingue ramal único de ramais de entrada eletricamente independentes. A existência de dois eletrodutos corresponde, neste arranjo, ao encaminhamento dos conjuntos do mesmo ramal. Considerando a seção individual de 150~mm², a resistência nominal de 300~daN corresponde preliminarmente à faixa até 185~mm² para ramal único. A expressão de engastamento produz 1,50~m para um poste de 9~m, como resultado de cálculo \cite[p.~1 e~7]{palatiummemorial2025}\cite[pranchas~001 e~005]{palatiumpranchas2025}\cite[seção~6, item~7.16.5 e Desenho~1]{enel1580}.

A verificação do poste depende dos esforços reais, das condições de instalação, da aplicação da nota~2 da Tabela~18 e da avaliação da Enel (IP-0004). Pelo mesmo motivo, a configuração do ramal aéreo depende da potência efetiva do transformador, da disponibilidade da rede e da avaliação da Enel (IP-0005). A seção dos condutores de entrada corresponde ao trecho subterrâneo do projeto. O paralelismo desse trecho será examinado na análise do ramal de entrada, considerando que a alínea~l do item~7.19.6 integra as prescrições para rede de distribuição subterrânea \cite[itens~7.16.2 e~7.19.6 e tabelas~17--18]{enel1580}.

A seção de 150~mm² do ramal de entrada atende à condição de incidência da alínea~k do item~7.16.2. No acervo examinado, não foi identificado termo específico que ateste a resistência mecânica requerida. Sua apresentação foi classificada como adequação potencialmente necessária de natureza documental (IP-0006). A eventual necessidade de substituição do poste depende da avaliação de IP-0004 \cite[p.~1 e~7]{palatiummemorial2025}\cite[item~7.16.2, alínea~k]{enel1580}.
